\begin{figure}[t]
\centering
% Shared color TikZ styles for the paper figures.
\colorlet{paperBlue}{blue!9!white}
\colorlet{paperBlueLine}{blue!55!black}
\colorlet{paperTeal}{teal!10!white}
\colorlet{paperTealLine}{teal!55!black}
\colorlet{paperAmber}{orange!15!white}
\colorlet{paperAmberLine}{orange!60!black}
\colorlet{paperRose}{red!9!white}
\colorlet{paperRoseLine}{red!55!black}
\colorlet{paperViolet}{violet!10!white}
\colorlet{paperVioletLine}{violet!55!black}
\colorlet{paperInk}{black!72}

\tikzset{
  paperfig/.style={
    font=\small,
    >=Stealth,
    line cap=round,
    line join=round
  },
  paperline/.style={
    -{Stealth[length=2.1mm,width=1.6mm]},
    line width=0.58pt,
    draw=paperInk
  },
  paperdash/.style={
    paperline,
    dashed,
    draw=black!48
  },
  paperbox/.style={
    draw=black!58,
    line width=0.58pt,
    rounded corners=2pt,
    fill=white,
    align=center,
    inner xsep=5pt,
    inner ysep=5pt,
    minimum height=8mm
  },
  paperdata/.style={
    paperbox,
    fill=paperBlue,
    draw=paperBlueLine
  },
  paperproc/.style={
    paperbox,
    fill=paperTeal,
    draw=paperTealLine
  },
  papermodel/.style={
    paperbox,
    fill=paperAmber,
    draw=paperAmberLine,
    line width=0.65pt
  },
  paperout/.style={
    paperbox,
    fill=paperRose,
    draw=paperRoseLine,
    line width=0.65pt,
    font=\small\bfseries
  },
  paperstate/.style={
    paperbox,
    fill=paperViolet,
    draw=paperVioletLine
  },
  papernote/.style={
    font=\scriptsize,
    align=center,
    text=black!70,
    fill=white,
    inner sep=1.4pt
  },
  paperaxis/.style={
    line width=0.45pt,
    draw=black!72
  },
  papergrid/.style={
    line width=0.25pt,
    draw=black!15
  },
  paperplot/.style={
    line width=0.95pt,
    draw=paperBlueLine
  },
  paperplotA/.style={
    line width=0.95pt,
    draw=paperBlueLine
  },
  paperplotB/.style={
    line width=0.95pt,
    draw=paperAmberLine
  },
  papermark/.style={
    circle,
    draw=paperBlueLine,
    fill=white,
    line width=0.55pt,
    inner sep=1.6pt
  },
  papermarkA/.style={
    papermark,
    draw=paperBlueLine,
    fill=paperBlue
  },
  papermarkB/.style={
    papermark,
    draw=paperAmberLine,
    fill=paperAmber
  }
}

\resizebox{0.98\linewidth}{!}{%
\begin{tikzpicture}[paperfig, x=1mm, y=1mm]
  \node[paperdata, text width=24mm] (task) at (0,32) {任务集合\\$\mathcal{D}=\{x_i\}$};
  \node[paperstate, text width=21mm] (budget) at (43,32) {预算档位\\$b$};
  \node[papermodel, text width=30mm] (orch) at (21.5,19) {编排器\\生成候选 $y_i^{(b)}$};
  \node[paperproc, text width=30mm] (exec) at (21.5,6) {统一执行器 $E$\\二值判定 $z_i^{(b)}$};
  \node[paperout, text width=20mm] (pass) at (0,-7) {通过率\\$p_b$};
  \node[paperout, text width=20mm] (calls) at (21.5,-7) {调用次数\\$c_i^{(b)}$};
  \node[paperout, text width=20mm] (tokens) at (43,-7) {令牌成本\\$t_i^{(b)}$};

  \draw[paperline] (task) -- (orch);
  \draw[paperline] (budget) -- (orch);
  \draw[paperline] (orch) -- (exec);
  \draw[paperline] (exec) -- (pass);
  \draw[paperline] (exec) -- (calls);
  \draw[paperline] (exec) -- (tokens);

  \node[papernote, text width=46mm] at (21.5,-17)
    {同一测试、超时与判定口径下比较收益和成本};
\end{tikzpicture}%
}
\caption{统一预算记账协议}
\label{fig:budget_accounting_protocol}
\end{figure}
